\chapter{Inverse, racine carrée et cube}
\section*{Exercice 15.1 -- Inverse}
La première donne $x=1/4$. Pour la seconde, étudier $(1+2x)/x\leq0$ : le numérateur s'annule en $-1/2$, le dénominateur en $0$ interdit. Le quotient est négatif ou nul sur $[-1/2,0[$.

\section*{Exercice 15.2 -- Racine}
Le domaine exige $x\geq1$. L'égalité donne $x-1=16$, soit $17$. L'inégalité donne $0\leq x-1<16$, donc $[1,17[$.

\section*{Exercice 15.3 -- Cube}
Le cube est strictement croissant sur $\R$. L'égalité donne $x=-3$ ; l'inégalité équivaut à $x<2$, soit $]-\infty,2[$.

\section*{Exercice 15.4 -- Comparer des images}
Sur $]-\infty,0[$, l'inverse décroît ; puisque $-4<-2$, $1/(-4)>1/(-2)$. La racine croît sur les positifs, donc $\sqrt{11}<\sqrt{12}$.

\section*{Exercice 15.5 -- Solution parasite}
Il faut $x+6\geq0$ et surtout $x\geq0$, puisque la racine est non négative. Sur $[0,+\infty[$, l'égalité équivaut à $x+6=x^2$, donc $(x-3)(x+2)=0$. Seul $3$ convient ; $-2$ est une solution de l'équation carrée, pas de l'initiale.

\section*{Exercice 15.6 -- Symétrie}
Un point du carré a la forme $(t,t^2)$ avec $t\geq0$. Son image est $(t^2,t)$ ; comme $\sqrt{t^2}=t$, il appartient à la courbe de la racine. Réciproquement, tout point $(x,\sqrt x)$ s'obtient avec $t=\sqrt x$.
